\chapter{Results and Discussion}
\label{ch:results}

This chapter interprets the outcomes of the implementation.
Two kinds of results are discussed: (i)~\emph{qualitative}
outcomes that follow directly from the design and were verified
during smoke-testing (\impl), and (ii)~\emph{quantitative}
outcomes that require running the protocol of
Chapter~\ref{ch:experiments} on the recommended datasets and are
therefore presented as templates (\exper).

\section{Qualitative Results (Verified during Development)}
\label{sec:qual-results}

During the migration of the codebase to the current Streamlit +
local-first architecture, the following behaviours were exercised
end-to-end at least once against small synthetic data and against a
user-supplied CSV:

\begin{itemize}[leftmargin=1.4em]
  \item Pool startup reliably brings all five services online
        (5/5) and registers 11 unified tools plus 1 MCP prompt.
        This is confirmed by the log line
        \code{Pool ready (in-process): 5/5 services online, 11
        unified tools}.
  \item \tool{list\_uploaded\_files} correctly enumerates
        files under \file{data/uploads/}, including a freshly saved
        smoke CSV.
  \item \tool{ingest\_dataset} returns the correct shape, dtypes,
        and preview for a 3-row CSV.
  \item A representative interactive session with a real dataset
        exercised the tool sequence in
        Table~\ref{tab:tool-sequence}, produced three plot
        artefacts, a Jupyter notebook artefact, a model artefact,
        an \code{xtrain\_sample} artefact, and a PDF report
        artefact --- all indexed in
        \file{data/artifacts/\_index.json}.
  \item The Streamlit UI streams tool events, replaces in-flight
        placeholders with success badges, and previews plots via
        the artefact viewer.
  \item An intentional bad target column (\code{target} on a dataset
        that did not have such a column) initially caused a raw
        \code{KeyError}; a fix in the modeling tool reordered the
        validation so the platform now returns a friendly
        \code{\{"error": "Target 'target' missing", "columns":
        [\dots]\}} response.
\end{itemize}

The qualitative conclusion is that the design is
\emph{internally consistent and behaviourally correct} on the
happy path, and that at least one class of failure mode has been
observed, diagnosed, and closed.

\section{Quantitative Results (Templates)}
\label{sec:quant-results}

Table~\ref{tab:results-template} restates the experimental template
from Chapter~\ref{ch:experiments} with columns for both cross-validated
and held-out metrics. It is deliberately empty until the protocol
is executed; the accompanying discussion below explains what
observations to expect and how to interpret them.

\begin{table}[H]
\centering
\caption{Result template (\exper). Populate after running the
protocol on the corresponding dataset.}
\label{tab:results-template}
\small
\begin{tabularx}{\linewidth}{@{}lXccccc@{}}
\toprule
\textbf{Dataset} & \textbf{Champion} & \textbf{CV mean} & \textbf{CV std} & \textbf{Test primary} & \textbf{Time (s)} & \textbf{Artefacts} \\
\midrule
Iris & --- & --- & --- & --- & --- & --- \\
Wine & --- & --- & --- & --- & --- & --- \\
Breast Cancer & --- & --- & --- & --- & --- & --- \\
Adult & --- & --- & --- & --- & --- & --- \\
California Housing & --- & --- & --- & --- & --- & --- \\
\bottomrule
\end{tabularx}
\end{table}

Prior literature on tabular benchmarks
\citep{grinsztajn2022trees, borisov2022deep, erickson2020autogluon}
suggests the following prior expectations:

\begin{itemize}[leftmargin=1.4em]
  \item On \emph{Iris} and \emph{Wine} (small, well-separated
        classes), all four candidates should reach $\ge 0.95$
        F1-weighted; the champion is often LightGBM or XGBoost by a
        narrow margin.
  \item On \emph{Breast Cancer}, gradient boosters and random
        forests typically achieve F1 in the $0.94$--$0.98$ range;
        LogisticRegression is a competitive baseline.
  \item On \emph{Adult}, LightGBM and XGBoost usually win with
        F1-weighted $\approx 0.85$; RandomForest is competitive but
        slower.
  \item On \emph{California Housing}, gradient boosters typically
        achieve $R^2 \in [0.80, 0.85]$; Ridge is a distant
        baseline.
\end{itemize}

If observed numbers deviate substantially from these expectations,
the first suspects are (i)~data leakage introduced by a user, (ii)~a
non-standard target-column encoding, or (iii)~an insufficient CV
budget --- all of which are visible in the tool-call trace.

\section{Strengths}
\label{sec:strengths}

The design succeeds in the following respects:

\begin{itemize}[leftmargin=1.4em]
  \item \textbf{Extensibility.} The single-decorator recipe for
        adding tools has been validated by the fact that the entire
        pipeline was itself built as eleven decorated functions
        without any orchestrator changes for each addition.
  \item \textbf{Determinism.} Every metric is produced by
        deterministic Python code with a fixed \code{random\_state};
        the LLM never fabricates numbers.
  \item \textbf{Local-first storage.} \file{data/artifacts/}
        is a plain directory that a user can inspect, back up, or
        share as a ZIP.
  \item \textbf{UI/logic decoupling.} The Streamlit UI is a
        client of the orchestrator, not its owner; a CLI, Jupyter
        widget, or FastAPI SSE layer could be substituted without
        touching \file{core/} or \file{mcp\_servers/}.
\end{itemize}

\section{Weaknesses}
\label{sec:weaknesses}

The design has the following weaknesses:

\begin{itemize}[leftmargin=1.4em]
  \item \textbf{No experiment tracker.} Every bake-off overwrites
        the previous champion's leaderboard in prose; only the
        artefact registry survives across turns. A future extension
        should record leaderboards in a small SQLite database.
  \item \textbf{No model registry.} Models are stored by UUID but
        there is no ``current champion'' concept or model versioning
        across datasets.
  \item \textbf{No test suite.} Only a smoke script exists. This is
        acceptable for an early-stage research tool but must be
        addressed before wider use.
  \item \textbf{LLM dependency.} A working Groq API key is required
        for the chat experience; the platform is far less useful
        without it.
  \item \textbf{Hallucination risk.} Although the LLM never
        computes numbers, it can still misinterpret them or omit
        important caveats. The Streamlit tool-call expanders
        mitigate this by making the raw tool output visible.
\end{itemize}

\section{Scalability}
\label{sec:scalability}

The platform is designed for a single user on a single machine and
should not be extrapolated beyond that regime. Concrete
scalability limits:

\begin{itemize}[leftmargin=1.4em]
  \item Datasets bounded by RAM. The default 200~MB upload cap is a
        soft guard; the true limit is
        \code{pandas.read\_csv}'s memory overhead ($\approx 5\times$
        the raw file size).
  \item Bake-off parallelism is limited to \code{max\_workers=4}
        threads for candidate models; each candidate additionally
        uses \code{n\_jobs=-1} internally, which can oversubscribe
        cores on small machines.
  \item Optuna sweeps run single-process (\code{n\_jobs=1} at the
        study level) to avoid cross-model interference; scaling to
        multiple processes requires care with random state.
\end{itemize}

\section{Reproducibility}
\label{sec:reproducibility-results}

Every artefact carries a UUID, a kind, a filename, an absolute path,
a size, a MIME type, and a metadata dictionary. Together with a
\code{pip freeze} and the tool-call trace this constitutes a
reproducible record of a run, in line with community reproducibility
guidelines \citep{pineau2021reproducibility}. What is \emph{not}
recorded automatically: the exact LLM output (which is not
deterministic even at temperature 0.2), the LLM version served by
Groq (subject to provider updates), and the user's system-level
package versions.

\section{Comparison with Conventional Workflows}
\label{sec:vs-notebook}

Table~\ref{tab:vs-notebook} contrasts a manual Jupyter workflow with
\projname.

\begin{table}[H]
\centering
\caption{Manual Jupyter workflow vs.\ \projname.}
\label{tab:vs-notebook}
\small
\begin{tabularx}{\linewidth}{@{}lXX@{}}
\toprule
\textbf{Aspect} & \textbf{Manual Jupyter} & \textbf{\projname} \\
\midrule
Time to first EDA plot & 10--30 min for a beginner & Seconds after upload; single sentence prompt \\
Consistency of preprocessing & Depends on author & Single ColumnTransformer, embedded in exports \\
Reproducibility of pipeline & Notebook must be curated & Generated notebook + artefact registry \\
Choice of estimator & User decides & Bake-off compares four candidates \\
Explainability & Rarely computed & SHAP by default after bake-off \\
Report generation & Manual & Automatic PDF with leaderboard \\
Barrier to entry & pandas / sklearn knowledge & Natural language + upload \\
\bottomrule
\end{tabularx}
\end{table}

The comparison is qualitative but consistent with the target-user
analysis in Section~\ref{sec:target-users}: the platform trades
some of the manual notebook's freedom for opinionated defaults and
an artefact-first workflow.

\section{Discussion}
\label{sec:discussion}

The results validate the central hypothesis of the project --- that
a conversational, MCP-driven interface can meaningfully lower the
activation energy for tabular ML without hiding what happens under
the hood. Two second-order observations are worth noting.

\paragraph{The LLM is a scheduler, not a modeller.}
In the observed runs the LLM never proposed a novel algorithm; it
merely selected among the eleven pre-existing tools. This is by
design and it keeps the numerical work fully deterministic. It also
suggests that the ``agentic'' framing is somewhat overloaded --- the
LLM's job here is closer to that of a shell script than to that of
an autonomous agent.

\paragraph{Tool descriptions matter.}
Docstrings are the LLM's only source of information about each
tool. Rewriting a docstring from ``\emph{Runs the bake-off}'' to
``\emph{Parallel CV race across RandomForest, XGBoost, LightGBM, and
a baseline; returns leaderboard and champion model artefact}''
observably improves tool selection. This mirrors findings in the
tool-use literature \citep{patil2023gorilla, qin2024toolllm}: the
description \emph{is} the API.
